\subsection{POLISH SPACES}

DEFINITION 1. A topological space X is said to be Polish if it is metrizable of countable type (§ 2, no. 8) and if there is a metric, compatible with the topology of X, with respect to which X is complete.

PROPOSITION I. a) Every closed subspace of a Polish space is Polish.

\begin{enumerate}
\def\labelenumi{\alph{enumi})}
\setcounter{enumi}{1}
\item
  The product of a countable family of Polish spaces is Polish.
\item
  The sum of a countable family of Polish spaces is Polish.
\end{enumerate}

Every subspace of a metrizable space of countable type is metrizable of countable type, and every closed subspace of a complete space is complete (Chapter II, § 3, no. 4, Proposition 8). Every countable product of metrizable spaces of countable type is again metrizable of countable type (§ 2, no. 8), and every countable product of complete metric spaces is a complete metric space with respect to a metric compatible with its topology (Chapter II, § 3, no. 5, Proposition 10 and Chapter IX, § 2, no. 4, Theorem 1, Corollary 2). Finally, let \((\mathbf{X}_n)\) be a sequence of non-empty Polish spaces, and consider the product space \(\mathbf{Y} = \mathbf{N} \times \prod_{n} \mathbf{X}_n\) , where \(\mathbf{N}\) carries the discrete topology; \(\mathbf{Y}\) is a Polish space by what has already been proved. On the other hand, let \(a_n\) be a point of \(\mathbf{X}_n\) for each \(n\) , and let \(f_n\) be the mapping of \(\mathbf{X}_n\) into \(\mathbf{Y}\) such that for each \(x \in \mathbf{X}_n\) we have

\[
f _ {n} (x) = (n, \left(y _ {p}\right)),
\]

where \(y_{p} = a_{p}\) if \(p \neq n\) and \(y_{n} = x\) . If \(X\) is the topological sum of the \(X_{n}\) , it is clear that the mapping \(f\) of \(X\) into \(Y\) which agrees with \(f_{n}\) on \(X_{n}\) for each \(n\) is a homeomorphism of \(X\) onto \(f(X)\) ; also, for each \(n\) , \(f_{n}(X_{n})\) is closed in \(Y\) , and the family \((f_{n}(X_{n}))\) is locally finite because \(N\) is discrete; therefore \(f(X) = \bigcup_{n} f_{n}(X_{n})\) is closed in \(Y\) (Chapter I, § 1, no. 5, Proposition 4), and thus \(f(X)\) is a Polish space, by a).

PROPOSITION 2. Every open subspace of a Polish space is Polish.

Let X be a Polish space, let d be a metric on X compatible with its topology, and let \(U \neq X\) be an open subset of X. Let V be the subset of \(R \times X\) consisting of all points \((t, x)\) such that

\[
t. d (x, \mathrm{X} - \mathrm{U}) = \mathrm{I};
\]

V is closed by Proposition 3 of § 2, no. 2 and is therefore Polish (Proposition 1). Since the restriction to V of the projection \(\mathbf{pr}_{2}:\mathbf{R}\times\mathbf{X}\to\mathbf{X}\) is a homeomorphism of V onto U (§ 2, no. 2, Proposition 3), U is a Polish subspace of X.

COROLLARY. Every locally compact, σ-compact, metrizable space X is Polish.

Let \(X'\) be the compact space obtained by adjoining a point at infinity to \(X\) ; \(X'\) is metrizable and of countable type (§ 2, no. 9, Corollary to Proposition 16), and \(X'\) is complete with respect to its unique uniformity (Chapter II, § 4, no. 1, Theorem 1). Hence \(X'\) is a Polish space; since \(X\) is open in \(X'\) , it follows that \(X\) is Polish.

PROPOSITION 3. Let X be a Hausdorff topological space. Then the intersection of a sequence \((\mathbf{A}_{n})\) of Polish subspaces of X is a Polish subspace.

Let \(f\) be the diagonal mapping of \(\mathbf{X}\) into \(\mathbf{X}^{\mathbf{N}}\) {[}Set Theory, Chapter II, § 5, no. 3; recall that \(f(x) = (y_n)\) where \(y_n = x\) for all \(n \in \mathbb{N}\) {]}. We shall use the following lemma:

LEMMA I. Let \((\mathbf{A}_{n})\) be a sequence of subsets of a Hausdorff topological space X. Then the restriction of the diagonal mapping \(f: X \to X^{N}\) to the subspace \(\bigcap_{n} A_{n}\) of X is a homeomorphism of \(\bigcap_{n} A_{n}\) onto a closed subspace of \(\prod_{n} A_{n}\) .

This image is the intersection of \(\prod_{n} A_n\) and the diagonal \(\Delta = f(X)\) , which is closed in \(X^N\) because \(X\) is Hausdorff (Chapter I, § 8, no. 1); and \(f\) is a homeomorphism of \(X\) onto \(\Delta\) .

With the hypotheses of Proposition 3, \(\prod_{n} A_{n}\) is a Polish space (Proposition 1), hence \(\bigcap_{n} A_{n}\) is Polish by Lemma 1 and Proposition 1.

COROLLARY. The set of irrational numbers, endowed with the topology induced by that of the real line R, is a Polish space.

It is the intersection of a countable family of open sets in R, namely the complements of sets consisting of a single rational number.

THEOREM 1. A subspace Y of a Polish space X is Polish if and only if Y is the intersection of a countable family of open sets in X.

The sufficiency of the condition follows immediately from Propositions 2 and 3. To show necessity, let \(d\) be a metric compatible with the topology of Y and with respect to which Y is complete. Let \(\overline{\mathbf{Y}}\) be the closure of Y in X. For each integer \(n>0\) , let \(\mathbf{Y}_{n}\) be the set of all \(x\in\overline{\mathbf{Y}}\) which have an open neighbourhood U such that the diameter of \(\mathbf{U}\cap\mathbf{Y}\) (with respect to the metric \(d\) ) is \(\leqslant1/n\) . Clearly \(\mathbf{Y}_{n}\) is open in \(\overline{\mathbf{Y}}\) and contains Y. Let \(x\in\bigcap_{n}\mathbf{Y}_{n}\) ; then \(x\in\overline{\mathbf{Y}}\) , and the trace on Y of the neighbourhood filter of \(x\) in X is a Cauchy filter (with respect to \(d\) ); hence this filter converges to a point of Y, and thus \(x\in\mathbf{Y}\) . Hence \(\mathbf{Y}=\bigcap_{n}\mathbf{Y}_{n}\) . For each \(n\) , let \(\mathbf{H}_{n}\) be an open subset of X such that \(\mathbf{H}_{n}\cap\overline{\mathbf{Y}}=\mathbf{Y}_{n}\) , and let \((\mathbf{U}_{m})\) be a sequence of open subsets of X such that \(\overline{\mathbf{Y}}=\bigcap_{m}\mathbf{U}_{m}\) ( \(\S\) 2, no. 5, Proposition 7); then Y is the intersection of the countable family of open sets \((\mathbf{H}_{n}\cap\mathbf{U}_{m})\) .

COROLLARY I. A space X is Polish if and only if it is homeomorphic to a countable intersection of open sets in the cube IN, where I is the interval {[}O, I{]} of R.

The condition is clearly sufficient, and it is necessary because every metrizable space of countable type is homeomorphic to a subspace of \(I^{N}\) ( \(\S\) 2, no. 8, Proposition 12).

COROLLARY 2. Let X and Y be two Polish spaces and let \(f: \mathbf{X} \to \mathbf{Y}\) be a continuous mapping. If Z is a Polish subspace of Y, then \(\overline{f}^{1}(Z)\) is a Polish subspace of X.

For \(Z = \bigcap_{n} Z_n\) , where the \(Z_n\) are open subsets of \(Y\) ; hence

\[
\overline {{{f}}} ^ {1} (Z) = \bigcap_ {n} \overline {{{f}}} ^ {1} (Z _ {n}),
\]

and the sets \(\overline{f}^1 (Z_n)\) are open in \(\mathbf{X}\) .

